\documentclass[11pt]{article}
\usepackage[T1]{fontenc}
\usepackage[margin=1in]{geometry}
\usepackage{amsmath,amssymb}
\usepackage[hidelinks]{hyperref}

\newcommand{\R}{\mathbb R}

\title{Literature-Implied Answer to a Vector-Valued Null Universal Differentiability Question}
\author{Run \texttt{fa\_banach\_001}, agent \texttt{agent\_lane\_15}}
\date{2026-06-28}

\begin{document}
\maketitle

\section*{Status}

This is a lightweight literature-implied answer packet.  It records that a
finite-dimensional vector-valued question mentioned in the introduction of
Doré--Maleva \cite{DM2011} is answered by later work of Preiss--Speight
\cite{PS2014}, together with the Alberti--Csörnyei--Preiss and
Csörnyei--Jones results cited there.  This is not a new proof produced by the
run.

\section*{Source Question}

Doré--Maleva \cite{DM2011} prove that every nonzero Banach space with
separable dual contains a closed bounded scalar universal differentiability set
of Hausdorff dimension \(1\).  In the introduction, they also discuss the
finite-dimensional vector-valued analogue.  They ask whether there are
Lebesgue null sets in \(\R^n\) containing a differentiability point of every
Lipschitz map
\[
  f:\R^n\to\R^m
\]
when \(n\ge m\ge2\).  The source says that the answer was then known only for
\(m=2\), that the case \(n=m=2\) is negative, and that no similar positive
results were known when the codomain dimension is at least \(3\).

\section*{Supporting Answer}

Preiss--Speight \cite{PS2014} state in their introduction that the converse to
Rademacher's theorem is valid if and only if \(m\ge n\).  Equivalently:

\begin{itemize}
\item If \(m\ge n\), every Lebesgue null \(N\subset\R^n\) is contained in the
non-differentiability set of some Lipschitz map \(\R^n\to\R^m\).  Thus no
Lebesgue null universal differentiability set exists in this range.
\item If \(n>m\), the converse fails.  Hence there is a Lebesgue null
\(N\subset\R^n\) such that every Lipschitz \(f:\R^n\to\R^m\) is differentiable
at some point of \(N\).
\end{itemize}

Theorem 1 of \cite{PS2014} proves the endpoint positive case \(m=n-1\):
for every \(n>1\), there is a Lebesgue null set
\[
  N\subset\R^n
\]
containing a differentiability point of every Lipschitz
\[
  f:\R^n\to\R^{n-1}.
\]
The same introduction says that, by a simple modification explained later in
the paper, for every \(n>m\) and every \(\tau>0\) there is such an \(N\) of
Hausdorff dimension at most \(m+\tau\).

\section*{Identification}

The Doré--Maleva question is exactly the finite-dimensional null universal
differentiability set question for vector-valued Lipschitz maps.  The
Preiss--Speight theorem and cited converse results settle the dichotomy:
positive precisely for \(n>m\), negative for \(m\ge n\).  Therefore the
question should not be treated as a live target for new proof work in this
run.

\section*{Search and Scope Notes}

The local run indexes had no exact packet for arXiv:1103.5094.  The local
source corpus was searched for the arXiv id, the source title, ``universal
differentiability set'', ``Lebesgue null'', and the \(n,m\) vector-valued
phrasing.  The relation to \cite{PS2014} was found in local parsed sources:
arXiv:1304.6916 explicitly says that the converse to Rademacher's theorem is
completely answered, and cites Doré--Maleva for the scalar \(m=1\) UDS
developments.  The supporting paper does not appear to be written as an
answer to this exact paragraph of \cite{DM2011}, so the packet is classified
as literature-implied rather than literature-already-answered.

\begin{thebibliography}{9}

\bibitem{DM2011}
M. Doré and O. Maleva,
\emph{A universal differentiability set in Banach spaces with separable dual},
arXiv:1103.5094; J. Funct. Anal. 261 (2011), 1674--1710.

\bibitem{PS2014}
D. Preiss and G. Speight,
\emph{Differentiability of Lipschitz Functions in Lebesgue Null Sets},
arXiv:1304.6916.

\end{thebibliography}

\end{document}
